% ==============================================================================
% INTRODUCCIÓN — Sección Nivel 1 (APA 7)
% Contenido copiado textualmente de texto_formatter/introduccion.txt
% ==============================================================================

\section{Introducción}\label{sec:introduccion}

La solicitud de becas es un elemento esencial para garantizar la estabilidad académica y económica de los estudiantes durante su formación universitaria. Sin embargo, el proceso actual de calificación y selección de los beneficiarios de la \enquote{Beca Comedor} en la UAGRM es manual, tedioso y susceptible a demoras. Por ello, este proyecto propone el desarrollo de un sistema de información integral impulsado por Inteligencia Artificial (IA) para la preselección automatizada.

Esta tecnología permitirá optimizar los tiempos de revisión y filtrar de manera eficiente aquellas postulaciones que presenten inconsistencias o que no cumplan con los requisitos mínimos establecidos para obtener el beneficio. Adicionalmente, la propuesta aborda la necesidad de mejorar la comunicación y el análisis de datos mediante un ecosistema digital conectado. Una aplicación móvil facilitará a los universitarios la postulación remota y la consulta transparente del menú diario. Simultáneamente, un portal web dotado de herramientas de Business Intelligence (BI) proveerá a la administración de tableros interactivos para monitorear indicadores clave del servicio, apoyando una toma de decisiones estratégica e informada.

El estudio se justifica por la necesidad imperante de modernizar y agilizar este proceso administrativo. Para estructurar este desarrollo, se aplica el Proceso Unificado de Desarrollo de Software (PUDS) como marco metodológico, junto con el estándar IEEE 830 para la especificación de requisitos de software.

El presente documento detalla todo el ciclo de vida del proyecto, desde el análisis de requerimientos hasta su implementación técnica. Asimismo, detalla el diseño arquitectónico, entrenamiento de IA, escalabilidad y los estándares de calidad.

Esta investigación no solo contribuye a mejorar la eficiencia operativa en la administración de la DUBSS, sino que establece un marco metodológico y tecnológico de referencia para futuros desarrollos que integren modelos de IA para la automatización de procesos administrativos y la toma de decisiones basada en datos dentro de la universidad.